% Generated from locale/tr/content/lambda-calculus/syntax/eta.tex; do not hand-edit.


\olfileid[tr]{lam}{syn}{eta}
\olsection{$\eta$-dönüşümü}

$\lambda$-terimleri üzerinde başka bir bağıntı daha vardır. \olref[fv]{sec}
içinde bir argüman alıp $f$'yi ona uygulayan $\lambd[x][(fx)]$ örneğini
kullandık. Başka bir deyişle bu, $f$ ile aynı fonksiyondur:
$\lambd[x][(fx)]N$ ile $fN$ ifadelerinin ikisi de $fN$'ye indirgenir. Bu
düşünceyi $\eta$-indirgeme ve $\eta$-genişletmeyle ifade ederiz.

\begin{defn}[$\eta$-daraltma, $\eredone$]
  \ollabel{defn:beredone}
  \emph{$\eta$-daraltma} ($\eredone$), terimler üzerinde aşağıdaki koşulu
  sağlayan en küçük bağdaşır bağıntıdır:
  \[
    \lambd[x][M x] \eredone M \text{, eğer } x \notin \FV{M}.
  \]
\end{defn}

\begin{defn}[$\beta\eta$-indirgeme, $\bered$] \ollabel{defn:bered}
  \emph{$\beta\eta$-indirgeme} ($\bered$), $\bredone$ ve $\eredone$
  bağıntılarını içeren, terimler üzerindeki en küçük yansımalı ve geçişli
  bağıntıdır; yani yansıma ve geçişlilik kurallarına ek olarak şu iki kuralı
  içerir:
  \begin{enumerate}
  \item $M \bredone N$ ise $M \bered N$. \ollabel{defn:bered3}
  \item $M \eredone N$ ise $M \bered N$. \ollabel{defn:bered4}
  \end{enumerate}

\end{defn}

\begin{defn}
  $\equal$ denklik bağıntısını şu $\eta$-dönüşümü kuralıyla genişletiriz:
  \[
  \lambd[x][f x] \equal f
  \]
  ve genişletilmiş bağıntıyı $\equal[\eta]$ ile gösteririz.
\end{defn}

$\eta$-denkliği, lambda terimlerinin kaplamsallığıyla ilişkili olduğu için
önemlidir:

\begin{defn}[Kaplamsallık]
  $\equal$ denklik bağıntısını şu (\ext) kuralıyla genişletiriz:
  \begin{center}
  $x \notin \FV{MN}$ olmak koşuluyla, $Mx \equal Nx$ ise $M \equal N$.
  \end{center}
  Genişletilmiş bağıntıyı $\equal[\ext]$ ile gösteririz.
\end{defn}

Kabaca bu kural, fonksiyon olarak görülen iki terim aynı argüman için aynı
biçimde davranıyorsa eşit sayılmaları gerektiğini söyler.

Şimdi $\eta$ kuralının tam olarak kaplamsallığı sağladığını ve bundan başka
bir şey eklemediğini ispatlayacağız.

\begin{thm}
  $M \equal[\ext] N$ ancak ve ancak $M \equal[\eta] N$.
\end{thm}

\begin{proof}
  Önce $\equal[\eta]$ bağıntısının kaplamsallık kuralı altında kapalı olduğunu
  ispatlarız. Başka bir deyişle \ext{} kuralı $\equal[\eta]$ bağıntısına yeni
  bir şey eklemez. Böylece $\equal[\eta]$, $\equal[\ext]$ bağıntısını içerir;
  $M \equal[\ext] N$ ise $M \equal[\eta] N$ olur.

  $\equal[\eta]$ bağıntısının \ext{} altında kapalı olduğunu göstermek için,
  \ext{} kuralıyla elde edilen herhangi bir $M\equal N$ sonucunun öncülü
  olarak $Mx\equal[\eta]Nx$ bulunduğuna dikkat edin. $\equal$ bağıntısının
  bir kuralı gereği
  $\lambd[x][Mx]\equal[\eta]\lambd[x][Nx]$ olur; iki tarafa $\eta$ kuralını
  uygulayınca $M\equal[\eta]N$ elde edilir.

  Benzer biçimde $\eta$ kuralının $\equal[\ext]$ içinde bulunduğunu
  ispatlarız. $x\notin\FV{M}$ olan herhangi bir $\lambd[x][Mx]$ ile $M$ için
  $(\lambd[x][Mx])x\equal[\ext]Mx$ vardır; \ext{} kuralı da
  $\lambd[x][Mx]\equal[\ext]M$ sonucunu verir.
\end{proof}

